% from ega0-16-fr.tex
\oldpage[0\textsubscript{IV-1}]{30}
$\dim(M/xM)=\dim(M)-1$, aucun des $\mathfrak{p}_i$ ne peut appartenir à
$\operatorname{Supp}(M/xM)$, et les idéaux premiers de ce support sont ceux
qui contiennent $\operatorname{Ann}(M)+Ax$~; comme les $\mathfrak{p}_i$
contiennent $\operatorname{Ann}(M)$, ils ne peuvent contenir $x$.
\end{proof}

\begin{corollary}[16.3.5]
\label{0.16.3.5-fr}
Avec les notations de (16.3.4), pour tout idéal
$\mathfrak{a}\subset\mathfrak{p}_i$, on a
$\dim(M/\mathfrak{a}M)=\dim(M)$ et si l'on pose
$N=M/\mathfrak{a}M$, la relation $\dim(M/xM)=\dim(M)-1$ entraîne
$\dim(N/xN)=\dim(N)-1$.
\end{corollary}

\begin{proof}
En effet, $\operatorname{Supp}(N)$ est l'ensemble des idéaux premiers
contenant $\mathfrak{a}+\operatorname{Ann}(M)$, donc $\mathfrak{p}_i$
appartient à $\operatorname{Supp}(N)$ et comme
$\dim(N)\leq\dim(M)=\dim(A/\mathfrak{p}_i)$, on a
$\dim(A/\mathfrak{p}_i)=\dim(N)$~; en outre les idéaux premiers
$\mathfrak{q}\in\operatorname{Supp}(N)$ tels que
$\dim(A/\mathfrak{q})=\dim(N)$ sont certains des $\mathfrak{p}_j$~; si $x$
n'appartient à aucun des $\mathfrak{p}_j$, on a donc
$\dim(N/xN)=\dim(N)-1$ en vertu de (16.3.4).
\end{proof}

\begin{definition}[16.3.6]
\label{0.16.3.6-fr}
Soient $A$ un anneau semi-local noethérien, $\mathfrak{r}$ son radical, $M$ un
$A$-module de type fini, et posons $n=\dim(M)$. On appelle \emph{système de
paramètres pour $M$} tout système $(x_i)_{1\leq i\leq n}$ de $n$ éléments de
$\mathfrak{r}$ tel que $M/(x_1M+\cdots+x_nM)$ soit de longueur finie.
\end{definition}

On a vu (16.2.3) qu'il existe toujours de tels systèmes.

On a vu au cours de la démonstration de (16.2.3.3) que si
$\mathfrak{J}=\sum_{i=1}^n x_iA$ est tel que $M/\mathfrak{J}M$ soit de
longueur finie, $\mathfrak{q}=\mathfrak{J}+\operatorname{Ann}(M)$ est un
\emph{idéal de définition} de $A$, et réciproquement, s'il en est ainsi,
$M/\mathfrak{J}M$ est un module de type fini sur l'anneau artinien
$A/\mathfrak{q}$, donc est de longueur finie. Il revient donc au même de dire
que $(x_i)_{1\leq i\leq n}$ est un système de paramètres \emph{pour $M$} ou
\emph{pour $A/\operatorname{Ann}(M)$} (ou encore que leurs images dans
$A/\operatorname{Ann}(M)$ forment un système de paramètres pour cet anneau).
Si $(x_i)_{1\leq i\leq n}$ est un système de paramètres pour $M$, il en est de
même de $(x_i^k)$ pour tout entier $k\geq1$, car on peut se borner, en vertu de
ce qui précède, au cas où $M=A$, et si $\mathfrak{J}$ est un idéal de
définition de $A$, l'idéal $\sum_{i=1}^n x_i^kA$ contient
$\mathfrak{J}^{kn}$, donc est aussi un idéal de définition.

\begin{proposition}[16.3.7]
\label{0.16.3.7-fr}
Soient $A$ un anneau semi-local noethérien, $\mathfrak{r}$ son radical,
$x_1,\ldots,x_k$ des éléments de $\mathfrak{r}$, $M$ un $A$-module de type
fini. On a alors
\begin{equation}
\label{0.16.3.7.1-fr}
\tag{16.3.7.1}
  \dim(M/(x_1M+\cdots+x_kM))\geq\dim(M)-k.
\end{equation}
En outre, pour que les deux membres de (16.3.7.1) soient égaux, il faut et il
suffit qu'il existe des éléments $x_i$
$(k+1\leq i\leq n=\dim(M))$ de $\mathfrak{r}$ tels que
$(x_i)_{1\leq i\leq n}$ soit un système de paramètres pour $M$.
\end{proposition}

\begin{proof}
L'inégalité (16.3.7.1) résulte de (16.3.4.1) par récurrence sur $k$. Si les
deux membres de (16.3.7.1) sont égaux, et si $x_{k+1},\ldots,x_n$ est un
système de paramètres pour $M/(x_1M+\cdots+x_kM)$, il est clair que
$M/(x_1M+\cdots+x_kM+x_{k+1}M+\cdots+x_nM)$ est de longueur finie, donc
$(x_i)_{1\leq i\leq n}$ est un système de paramètres pour $M$~;
réciproquement, s'il existe des $x_i$ $(k+1\leq i\leq n)$ ayant cette
propriété et si l'on pose $N=M/(x_1M+\cdots+x_kM)$, il est clair que
$N/(x_{k+1}N+\cdots+x_nN)$ est de longueur finie, donc
$\dim(N)\leq n-k$, et les deux membres sont égaux en vertu de (16.3.7.1).
\end{proof}

